\chapter{Background and Problem Statement}\label{ch:background}

This chapter fixes the security requirements, the threat model and the building blocks the scheme
uses, and ends with the problem statement.

\section{Security Requirements}

\begin{enumerate}
  \item \textbf{Authentication and integrity:} a receiver accepts a message only if it comes from a
  vehicle registered by the TA and has not been altered.
  \item \textbf{Anonymity:} neither RSUs nor other vehicles learn a vehicle's real identity from
  authentication or safety messages.
  \item \textbf{Unlinkability:} an observer should not be able to link the pseudonyms of one vehicle
  over time; at least, changing pseudonyms must be possible and effective.
  \item \textbf{Conditional traceability:} an authorized investigation can reveal the identity behind
  a pseudonym, but no single party can do so alone.
  \item \textbf{Revocation:} a misbehaving pseudonym, and after an investigation the whole vehicle,
  can be excluded quickly from the network.
  \item \textbf{Authorization:} an authenticated vehicle may still be denied a service its class is
  not entitled to (for example, a private car asking for emergency priority).
  \item \textbf{Replay resistance:} an old proof or message cannot be accepted again.
  \item \textbf{Performance:} authentication within a few milliseconds and safety-message overhead
  small enough to keep the shared channel uncongested.
\end{enumerate}

\section{Threat Model}\label{sec:threat}

\begin{itemize}
  \item \textbf{Outsider attackers} can listen to every message (a global passive eavesdropper),
  inject, replay and modify messages, but hold no registered credential.
  \item \textbf{Insider attackers} are registered vehicles that send false information: false
  positions or speeds, sudden-stop claims, message floods, and, in the emergency extension, false
  emergency claims. They may use all of their own pseudonyms.
  \item \textbf{RSUs} follow the protocol but are curious: they try to learn identities. Up to $f$ of
  the $n = 3f+1$ RSU organizations may be compromised and report honest vehicles; the trust-weighted
  extension (Section~\ref{sec:trust}) addresses this case.
  \item \textbf{The TA} registers vehicles and is trusted for registration. It cannot unmask a
  pseudonym alone: identity records are escrowed under a key shared among the RSU organizations.
  \item \textbf{The ledger} is a permissioned blockchain whose ordering service tolerates crash
  faults (Raft) or Byzantine faults (SmartBFT).
\end{itemize}
Physical attacks on OBUs, radio jamming and denial of service by exhausting the radio channel are
outside the scope.

\section{Cryptographic Building Blocks}

All public-key operations use the elliptic curve secp256k1~\cite{sec2} with generator $G$ of prime
order $q$, and the hash function $H$ is SHA-256~\cite{fips1804}.

\subsection{Schnorr Proof of Knowledge}

A vehicle holds a secret key $x \in \mathbb{Z}_q$ for each pseudonym, with public key $Y = xG$. The
Schnorr protocol~\cite{schnorr1991}, made non-interactive with the Fiat--Shamir
transform~\cite{fiat1986}, proves knowledge of $x$ without revealing it:
\begin{align}
  \text{commit:}\quad & r \xleftarrow{\$} \mathbb{Z}_q,\quad T = rG \\
  \text{challenge:}\quad & c = H(G \,\|\, Y \,\|\, T \,\|\, \mathit{ctx}) \\
  \text{response:}\quad & s = r + c\,x \bmod q \\
  \text{verify:}\quad & sG \stackrel{?}{=} T + cY
\end{align}
The context $\mathit{ctx}$ binds the proof to one RSU, one fresh RSU nonce and timestamp, the
requested service, the pseudonym and the vehicle class, so a proof cannot be replayed at another RSU,
later, or for another service. The proof is \emph{sound} (without $x$ a forger succeeds with
probability about $1/q$) and \emph{zero-knowledge} (a simulator that picks $c, s$ first and sets
$T = sG - cY$ produces proofs with the same distribution, so a real proof reveals nothing about $x$).
The commitment $T$ does not depend on the RSU, so a vehicle can precompute a pool of commitments
while idle and only compute $s$ when it meets an RSU.

\subsection{Merkle Trees}

A Merkle tree~\cite{merkle1987} hashes a list of leaves pairwise up to a single root. A path of
$\lceil \log_2 N \rceil$ sibling hashes proves that a leaf is in the tree. The TA builds one tree per
epoch over the leaves
\[
  \ell = \mathsf{enc}(\texttt{"VANET-LEAF-v1"},\ \mathit{pid},\ Y,\ \mathit{class}),
\]
where $\mathsf{enc}$ length-prefixes each field. Only the 32-byte root goes on the ledger. Because
the pseudonym, its key and the vehicle class are inside the leaf, a vehicle cannot claim another
pseudonym or class than the TA registered: the path would not verify.

\subsection{Threshold Identity Escrow}

The TA encrypts each vehicle's identity record (its real identity and its list of pseudonyms) with
AES-256 in GCM mode and splits the key with Shamir's $(t,n)$ secret sharing~\cite{shamir1979} among the $n$
RSU organizations, with $t = 2f+1$. An RSU releases its share only for a trace request that the ledger
shows as approved. Any $t$ shares reconstruct the key by Lagrange interpolation; fewer reveal nothing.

\subsection{Signatures, TESLA and Implicit Certificates}\label{sec:tesla-ecqv}

Tickets and handover requests use Schnorr signatures. The compact wire format uses
BIP-340~\cite{bip340} 64-byte signatures with 32-byte x-only keys.

TESLA~\cite{perrig2000tesla} replaces a signature per message with a message authentication code
(MAC) under a key from a one-way hash chain~\cite{lamport1981password} $K_i = H(K_{i+1})$. The sender uses $K_i$ during
interval $i$ and discloses it $d$ intervals later. A receiver buffers the message, checks the
disclosed key against the signed chain commitment, and then the MAC. A MAC is 16 bytes and about ten
times cheaper to compute than a signature; the cost is a delay of about $d$ intervals before a message
can be accepted.

An ECQV implicit certificate~\cite{sec4ecqv} replaces the RSU's signature on a ticket: the
certificate carries a point $P_U$ from which any receiver reconstructs the holder's public key as
$Q_U = H(\mathit{cert})\,P_U + Q_{CA}$. A certificate is about 64 bytes smaller than an explicit
key plus signature, and it is authenticated implicitly by the first signature verified under $Q_U$.

\subsection{BBS Anonymous Credentials}

With BBS signatures~\cite{boneh2004bbs,tessaro2023bbs} the TA signs a hidden vehicle secret, the
class and the epoch once. The vehicle then proves in zero knowledge that it holds such a signature,
and derives its pseudonym as $\mathit{nym}_j = H(\mathit{epoch}, j)^{sk_v}$ for $j < J$. Two proofs
cannot be linked to each other, and the scope $j$ limits the number of pseudonyms per epoch, which
bounds Sybil attacks. No per-pseudonym registration is needed.

\subsection{Bloom Filters}

A Bloom filter~\cite{bloom1970} represents a set in $m$ bits with $k$ hash functions. It has no false
negatives and a tunable false-positive rate. The extension uses one to put the recently revoked
pseudonyms into RSU beacons.

\section{Permissioned Blockchain: Hyperledger Fabric}

Hyperledger Fabric~\cite{androulaki2018fabric} is a permissioned blockchain in which every
participant has a certificate from its organization's Membership Service Provider (MSP). A
transaction is executed by the endorsing peers (\emph{execute}), ordered into blocks by the ordering
service (\emph{order}) and validated and committed by every peer (\emph{validate}). Smart contracts,
called chaincode, can read the caller's identity and therefore enforce role-based rules: only the TA
may register roots or policies, and only RSUs may record authentications or report misbehaviour.
The ordering service runs either Raft~\cite{ongaro2014raft}, which tolerates crashes of a minority of
orderers, or SmartBFT~\cite{barger2021smartbft}, which tolerates $f$ Byzantine orderers out of
$n \ge 3f+1$~\cite{lamport1982byzantine,castro1999pbft} and commits after $2f+1$ matching votes. Blocks are cut when a batch is full or after a
batch timeout, which is the main term in commit latency.

\section{Misbehaviour Detection}

Cryptography proves who sent a message, not that its content is true: a registered vehicle can still
broadcast a false position. Misbehaviour detection checks plausibility: a vehicle cannot move faster
than physically possible, accelerate beyond a limit, or be heard from beyond radio range. The VeReMi
dataset~\cite{heijden2018veremi} and its extension~\cite{kamel2020veremiext} are the standard
benchmark: realistic traces with 19 labelled attack types, among them constant or random position
and speed offsets, sudden stops, replay, denial of service and Sybil variants.

\section{Problem Statement}

Design a VANET authentication and authorization scheme that
\begin{enumerate}
  \item proves registration and vehicle class with a zero-knowledge proof bound to the RSU, the time,
  and the service, so that verifiers learn neither the identity nor the secret;
  \item keeps authorization, revocation and identity-reveal approvals on a permissioned blockchain
  without placing the blockchain on the path of individual safety messages;
  \item removes misbehaving vehicles with conditional traceability that no single party controls;
\end{enumerate}
and evaluate on a realistic 802.11p network with real traffic whether its latency, channel load and
detection performance are acceptable compared with certificate-based and centralized baselines.
